%\section{Economic Scenario Generator}
\label{Appendix: Stock Scenario Generator}

The method and algorithm presented in this appendix are used for the pathwise simulation of stock prices in Chapter \ref{section: Multi-dimensional Dynamic Portfolio Allocation}. We closely follow the approach presented in \cite{LUCIANO20101937, Luciano02072016}.

\section{Variance-Gamma processes and subordination}
Let $S_i=\{S_i(t)\}_{t \geq 0}$ denote a one-dimensional price process of stock $i$, for $i=1,\ldots,d$, on a filtered probability space $(\Omega, \F, \{\F\}_{t \geq 0}, \Prob)$. Suppose that the initial condition $S_i(0^-) = s_i$ is known and that $S_i$ follows a pure-jump SDE
\begin{equation}
    \diff S_i(t) = S_i(t^-) \left( \mu_i \diff t + \int_\R\left(e^z-1\right) \tilde N_i(\diff t, \diff z) \right),
    \label{D: SDE}
\end{equation}
where $\mu_i$ and $\tilde N_i$ are the drift coefficient and the compensated Poisson random measure of stock $i$, respectively. $\tilde N$ is given by
\begin{equation}
    \tilde N(\diff t, \diff z) = 
    \begin{cases}
        N(\diff t, \diff z) - \nu(\diff z) \diff t, & \text{for } |z| < 1 \\[1ex]
        N(\diff t, \diff z), & \text{for } |z| < 1.
    \end{cases}
    \label{D: N-tilde}
\end{equation}
In the stock price model \eqref{D: SDE}, the mark $z$ is understood to be the log-jump size. $\nu$ in \eqref{D: N-tilde} denotes the Lévy measure. For a Variance-Gamma (VG) process, it takes the form
\begin{equation*}
    \nu(\diff z) = \frac{C}{|z|} \left( e^{Gz} \ind{z<0} + e^{-Mz} \ind{z>0} \right)\diff z,
\end{equation*}
where $C,G,M > 0$ are constant parameters. We assume that every stock in the portfolio is a VG process with correlated parameters $C_i,G_i,M_i>0$. 

For a VG process, the SDE \eqref{D: SDE} has the solution
\begin{align}
    S_i(t) = s_i e^{(\mu_i+\omega_i)t + L(t;\sigma_i, \theta_i, \nu_i)}, \qquad S_i(0^-) = s_i,
    \label{D: Stock price model}
\end{align}
where 
\begin{equation}
    \omega_i = \frac{1}{\nu_i} \log \left( 1 - \theta_i\nu_i - \frac{1}{2}\sigma_i^2 \nu_i \right),
    \label{D: omega-i}
\end{equation}
and the log-return process $L=\{L(t)\}_{t \geq 0}$ is given by
\begin{align}
    L(t;\sigma_i, \theta_i, \nu_i) = \sigma_i W_i(G_i(t)) + \theta_iG_i(t), \qquad L(0^-)=0.
    \label{D: L-VG}
\end{align}
Here, $W_i$ denotes the Brownian motion of stock $i$ and $G_i$ is the corresponding subordinator. In particular, $G_i$ follows a Gamma distribution with shape $1/\nu_i$ and rate $\nu_i$, i.e.
\begin{equation*}
    G_i(t) \sim \text{Gamma}\left(\frac{t}{\nu_i}, \frac{1}{\nu_i} \right). 
\end{equation*}
The subordinator $G_i$ can be interpreted as a stochastic ``business clock'', enabling the stock price model to capture volatility clustering and non-normal returns. The clock accelerates during periods of high market activity, such as earning announcements or switches in market regime, and slows down during quieter periods. Moreover, the weighted jump term $\theta_i G_i$ in \eqref{D: L-VG} represents the arrival of information or large trades which lead to discontinuous jumps in the asset prices.

The marginal parameters $\sigma_i, \nu_i > 0$ and $\theta_i \in \R$ are related to $C_i$, $G_i$, and $M_i$ in the sense that
\begin{align*}
    & C_i = \frac{1}{\nu_i} > 0, \\
    & G_i = \frac{1}{\sqrt{\frac{1}{4} \theta_i^2 \nu_i^2 + \frac{1}{2}\sigma_i^2\nu_i} - \frac{1}{2}\theta_i\nu_i} > 0, \\
    & M_i = \frac{1}{\sqrt{\frac{1}{4} \theta_i^2 \nu_i^2 + \frac{1}{2}\sigma_i^2\nu_i} + \frac{1}{2}\theta_i\nu_i} > 0,
\end{align*}
see \cite[Chapter 5.3.7]{schoutens2003levy}. In practice, it is often preferable to estimate $\sigma_i$, $\theta_i$, and $\nu_i$ directly from market data, such as European call options, rather than working with $C_i$, $G_i$, and $M_i$, see Appendix \ref{appendix: Calibration Procedure}.

\section{Simulation of correlated stocks via multivariate subordination}
\subsection{Multivariate subordination}
Let $\mathbf{S}=\{(S_1(t), \ldots, S_d(t))\}_{t \geq 0}$ denote a multidimensional stock price process with $d \in [1,\infty)$, where each stock $S_i$ follows the VG process \eqref{D: Stock price model} with parameters $\sigma_i$, $\theta_i$, and $\nu_i$ for $i=1,\ldots,d$. Furthermore, define the multidimensional log-return process $\mathbf{L}=\{(L_1(t), \ldots, L_d(t))\}_{t \geq 0}$, where each $L_i$ follows \eqref{D: L-VG} with respective parameters for $i=1,\ldots,d$. 

Suppose that $Y_i = \{Y_i(t)\}_{t \geq 0}$, $i=1, \ldots,d$, and $Z = \{Z(t)\}_{t \geq 0}$ are independent subordinators, with
\begin{equation}
    Y_i(t) \sim \text{Gamma}\left( \frac{t}{\nu_i}-a, \frac{1}{\nu_i} \right) \quad \text{and} \quad Z(t) \sim \text{Gamma}\left( a,1 \right).
    \label{D: Y and Z samples}
\end{equation}
The parameter $a > 0$ in \eqref{D: Y and Z samples} is a real number. In particular, we choose
\begin{align*}
    a = \lambda \min \left\{ \frac{1}{\nu_1}, \ldots, \frac{1}{\nu_d} \right\},
\end{align*}
with $\lambda \in (0,1)$. Because $a$ is present in the shape of both $Y_i$ and $Z$, it is understood to be the common activity parameter. In the ``business clock'' interpretation of the subordinator, $Y_i$ is interpreted as the idiosyncratic clock of stock $i$, while $Z$ is the shared market clock. Thus, both the Gamma clocks and the jump components of the stocks in the portfolio are correlated through the parameter $a$. The choice of $\lambda$ depends on the data, but should not be too big as it would restrict the jumps of the risky asset with the smallest jump intensity $1/\nu$. In our case, we have chosen $\lambda = 0.4$ arbitrarily.

With the subordination processes $Y_i$ and $Z$ in mind, the Lévy process $\mathbf{G} = \{\mathbf{G}(t)\}_{t \geq 0}$ given by
\begin{equation*}
    \mathbf{G}(t) = \big(Y_1(t) + \nu_1 Z(t), \ldots, Y_d(t) + \nu_d Z(t)\big)
\end{equation*}
is a multivariate subordinator, see \cite{LUCIANO20101937}. 

Finally, let $W_i=\{W_i(t)\}_{t \geq 0}$, $i=1,\ldots, d$, a standard Brownian motion. We collect all Brownian motions into $\mathbf{W}(t)=(W_1(t), \ldots, W_d(t))$. Assuming that the stocks' Brownian motions are correlated through the matrix
\begin{equation}
    \mathbf{C} =
    \begin{pmatrix}
        1 & \rho_{12} & \cdots & \rho_{1d} \\
        \rho_{12} & 1 & \cdots & \rho_{2d} \\
        \vdots & \vdots & \ddots & \vdots \\
        \rho_{1d} & \rho_{2d} & \cdots & 1
    \end{pmatrix},
    \label{D: C-matrix}
\end{equation}
where $\rho_{ij} \in (-1,1)$ for all $i,j = 1, \ldots, d$ and $i\neq j$, we use Cholesky's decomposition theorem and write $\mathbf{C}=\mathbf{\Lambda \Lambda}^\top$. The diffusion matrix $\mathbf{A}$ is defined as
\begin{equation*}
    \mathbf A = \mathbf{D}\mathbf{\Lambda},
\end{equation*}
where $\mathbf{D} = \text{diag}\{\sigma_1, \ldots, \sigma_d\}$. Then, the log return process $\mathbf{L}$ follows the subordinated process
\begin{align}
    \mathbf{L}(t) = \mathbf{AW}(\mathbf{G}(t)) + \mathbf{\theta}\mathbf{G}(t), \qquad\mathbf{L}(0^-) = 0,
    \label{D: L-multidim}
\end{align}
where $\theta=(\theta_1, \ldots, \theta_d)$. The stock price process $\mathbf{S}$ is therefore
\begin{align}
    \mathbf{S}(t) = \mathbf{s} e^{(\mu + \omega)t + \mathbf{L}(t)}, \qquad\mathbf{S}(0^-) = \mathbf s
\end{align}
where $\mathbf{s}=(s_1, \ldots, s_d)$ is the initial vector, $\mu = (\mu_1, \ldots, \mu_d)$ is the drift vector, and $\omega=(\omega_1, \ldots, \omega_d)$. 

\subsection{Estimation of the Brownian correlation matrix}
The correlation matrix $\mathbf{C}$ cannot be directly observed from stock market data alone and must be estimated. In view of this, let $\mathbf{R}$ denote the correlation matrix of log-returns, which can be estimated using historical stock data. Following \cite{LUCIANO20101937, Luciano02072016}, suppose $\mathbf R$ takes the form
\begin{equation}
    \mathbf{R} =
    \begin{pmatrix}
        1 & \rho_{\mathbf{Y}}(1,2) & \cdots & \rho_{\mathbf{Y}}(1,d) \\
        \rho_{\mathbf{Y}}(1,2) & 1 & \cdots & \rho_{\mathbf{Y}}(2,d) \\
        \vdots & \vdots & \ddots & \vdots \\
        \rho_{\mathbf{Y}}(1,d) & \rho_{\mathbf{Y}}(2,d) & \cdots & 1
    \end{pmatrix},
    \label{D: R-matrix}
\end{equation}
where, for a Brownian motion correlation $\rho_{ij}$ in $\mathbf C$, $\rho_{\mathbf{Y}}(i,j)$ is given by
\begin{equation}
    \rho_{\mathbf{Y}}(i,j) = \frac{a \left(\rho_{ij} \sigma_i \sigma_j \sqrt{\nu_i \nu_j} + \theta_i \theta_j \nu_i \nu_j\right)}{\sqrt{(\sigma_i^2 + \theta_i^2 \nu_i)(\sigma_j^2 + \theta_j^2 \nu_j)}}.
    \label{D: rho-model}
\end{equation}

Then, if $ \rho^{\text{Market}}_{\mathbf{Y}}(i,j)$ denotes the observed log-return correlation between assets $i, j$ and $\rho^{\text{Model}}_{\mathbf{Y}}(i,j)$ is given by \eqref{D: rho-model}, we estimate $\rho_{ij}$ via constrained minimisation. Specifically, we solve
\begin{align*}
    & \min_{\substack{\hat \rho_{ij} \\ i,j=1,\ldots, d}}\sum_{i<j}\left( \rho^{\text{Market}}_\mathbf{Y}(i,j) - \rho^{\text{Model}}_\mathbf{Y}(i,j) \right)^2 \quad  \text{s.t } \hat \rho_{ij} \in (-1,1) \; \text{for all } i,j = 1,\ldots, d, \, i\neq j.
\end{align*}
The correlation matrix $\hat{\mathbf{C}}$ is then constructed using the estimated correlation coefficients $\hat \rho_{ij}$. 

\subsection{Numerical scheme}
\label{appendix: Numerical scheme}
With \eqref{D: L-multidim}-\eqref{D: rho-model}, the simulation of $d$-assets present in the portfolio is straightforward. Assume that we work on a uniform time grid $0=t_0 < t_1 < \cdots < t_{N_t}=T$ with $\Delta t = T/N_t$ and $N_t>2$. Furthermore, assume that the parameters $a$, $\mu=(\mu_1, \ldots, \mu_d)$, $\sigma=(\sigma_1, \ldots, \sigma_d)$, $\theta=(\theta_1, \ldots, \theta_d)$, $\mathbf{A}$, and $\hat{\mathbf{C}}$ are given.

At each time step $t_n$, we draw $N_{\text{MC}}$ (Monte Carlo samples) of $\mathbf{W}(t_n)$, $Y_i(t_n)$, $i=1,\ldots, d$, and $Z(t_n)$ according their respective distributions. Then all $d$ stock prices are computed simultaneously with
\begin{equation}
    \mathbf S(t_{n}) = \mathbf S(t_{n-1})e^{(\mu + \omega) \Delta t + \Delta \mathbf L(t_n)},
    \label{D: simulation of d stocks}
\end{equation}
where 
\begin{equation*}
    \Delta \mathbf L(t_n) = \mathbf{L}(t_{n}) - \mathbf{L}(t_{n-1}). 
\end{equation*}
The process is repeated for $n=1, \ldots, N_t$. Algorithm~\ref{alg: Economic Scenario Generator} summarises this simulation procedure. 
\begin{algorithm}[H]
\caption{Simulation procedure of $d$-correlated stocks via multivariate subordination}
\label{alg: Economic Scenario Generator}
\begin{algorithmic}[1]
\REQUIRE Parameters $a$, $\mu$, $\sigma$, $\nu$, $\theta$, $\mathbf{A}$, $\mathbf{C}$
\STATE Compute $\omega_i$ using \eqref{D: omega-i} for $i=1, \ldots, d$
\STATE Initialise time grid $0=t_0 < t_1 < \cdots < t_{N_t}$, with $\Delta t = T/N_t$ and $t_n-t_{n-1} = \Delta t$ for every $n=1,\ldots, N_t$
\STATE Initialise $\mathbf{S}_0=\mathbf{s}$
\FOR{$n = 1$ to $N_{t}$}
    \STATE Sample $N_{\text{MC}}$
    \begin{equation*}
        \Delta Y_{i,n} \sim \text{Gamma}\left(\frac{\Delta t}{\nu_i}, \frac{1}{\nu_i}\right), \quad \Delta Z_n \sim \text{Gamma} \left(a \Delta t,1\right), \quad \varepsilon_n \sim \mathcal{N}(0,\mathbf{I}_d)
    \end{equation*}
    \STATE Construct $\Delta \mathbf{G}_n = (\Delta Y_{1,n} + \Delta Z_n, \ldots, \Delta Y_{d,n} + \Delta Z_n)$
    \STATE Compute $\Delta \mathbf L_n = \mathbf A \sqrt{\Delta \mathbf{G}_n} \, \odot \varepsilon_n \Delta t + \theta \odot\Delta \mathbf{G}_n$
    \STATE Compute $\mathbf{S}_n = \mathbf{S}_{n-1}e^{\Delta \mathbf L_n}$
\ENDFOR
\end{algorithmic}
\end{algorithm}